% Generated draft: requires mathematical and source review.
\documentclass[11pt]{article}
\usepackage[margin=25mm]{geometry}
\usepackage{fontspec}
\setmainfont{Latin Modern Roman}
\usepackage{amsmath,amssymb,mathtools,unicode-math}
\setmathfont{Latin Modern Math}
\usepackage{enumitem,needspace,xurl,hyperref}
\setlength{\parindent}{0pt}
\setlength{\parskip}{5pt}
\setlength{\emergencystretch}{3em}
\newcommand{\N}{\mathbb{N}}
\newcommand{\Z}{\mathbb{Z}}
\newcommand{\Q}{\mathbb{Q}}
\newcommand{\R}{\mathbb{R}}
\newcommand{\C}{\mathbb{C}}
\newcommand{\F}{\mathbb{F}}
\newcommand{\A}{\mathbb{A}}
\newcommand{\PP}{\mathbb{P}}
\newcommand{\E}{\mathbb{E}}
\newcommand{\eps}{\varepsilon}
\newcommand{\dd}{\,\mathrm{d}}
\begin{document}
\section*{UnsolvedMath ID: 700170}
\textbf{Source problem code:} catalogue entry
\subsection{Definitions and mathematical statement}
For a group \(H\), define its derived series by \(H^{(0)}=H\) and \(H^{(k+1)}=[H^{(k)},H^{(k)}]\). The variety of soluble groups of derived length at most \(n\) consists of the groups satisfying \(H^{(n)}=1\).

Let \(X\) be a finite set, let \(F(X)\) be the free group on \(X\), and let \(r\in F(X)\). A group with the single defining relation \(r=1\) relative to this variety has the presentation
\[
G_{n,X,r}=F(X)\big/\bigl(F(X)^{(n)}\langle\!\langle r\rangle\!\rangle_{F(X)}\bigr),
\]
where \(\langle\!\langle r\rangle\!\rangle_{F(X)}\) denotes the normal closure of \(r\) in \(F(X)\). Its word problem is decidable if there is an algorithm that, on every input word \(w\) in \(X\cup X^{-1}\), halts and determines whether \(w=1\) in \(G_{n,X,r}\).

\textbf{Problem, under the finite generating set interpretation.} Is the word problem decidable for every \(G_{n,X,r}\), for every integer \(n\geq 3\), every finite generating set \(X\), and every relator \(r\in F(X)\)? Equivalently, for each such presentation, is there an algorithm deciding whether
\[
w\in F(X)^{(n)}\langle\!\langle r\rangle\!\rangle_{F(X)}
\]
for every input word \(w\)?

The supplied wording does not specify that the generating set is finite, nor whether an algorithm must be uniform over the presentations or over \(n\). Thus the formulation above is an explicitly qualified interpretation. A uniform version would require a single algorithm taking \((n,X,r,w)\) as input and deciding the same membership condition. The variety condition means derived length at most \(n\); it does not require that the resulting quotient have derived length exactly \(n\).
\subsection{Short English statement}
Kargapolov asks about the word problem for groups defined by one relation within the variety of soluble groups of derived length at most \(n\), where \(n\geq 3\). The draft states the finite generating set interpretation; the permitted generating sets and the intended uniformity of the algorithm remain unspecified in the supplied source.
\subsection{Sources}
\begin{itemize}
\item Local UnsolvedMath registry, supplied record with permanent ID \(700170\), exact source problem code \texttt{catalogue entry}, title \emph{Kourovka Notebook Problem 3.16}; upstream revision \texttt{d371ee917ac9647fbb06cd6fc529deec6f3dda82}.
\item Attributed in the supplied record to M.\ I.\ Kargapolov, proposed in 1969. The supplied background identifies \emph{The Kourovka Notebook (Unsolved Problems in Group Theory)}, 3rd issue (1969), Problem 3.16, and the 21st issue edition (2026), \texttt{arXiv:1401.0300v48}.
\item The supplied collection description identifies E.\ I.\ Khukhro and V.\ D.\ Mazurov as editors of the collection covering Issues 1--21 (1965--2026). The registry labels this problem \emph{open}; that label is supplied source metadata, not an independently verified literature assessment.
\end{itemize}
\end{document}
